\documentclass[aspectratio=169,11pt]{beamer}
\usetheme[progressbar=frametitle,numbering=fraction]{metropolis}

\usepackage{fontspec}
\usepackage[english,russian]{babel}
\usepackage{amsmath}
\usepackage{unicode-math}
\setsansfont{Calibri}
\setmainfont{Calibri}
\setmonofont{Consolas}[Scale=0.88]
\setmathfont{Cambria Math}
\usepackage{booktabs}
\usepackage{xcolor}

\definecolor{accent}{HTML}{2F6DB5}
\definecolor{good}{HTML}{2E7D4F}
\definecolor{bad}{HTML}{B23A48}
\setbeamercolor{frametitle}{bg=accent!85!black}
\setbeamercolor{progress bar}{fg=accent}
\setbeamercolor{alerted text}{fg=bad}
\newcommand{\pos}[1]{\textcolor{good}{#1}}
\newcommand{\negv}[1]{\textcolor{bad}{#1}}
\newcommand{\fc}{\mathrm{fc}}
\DeclareMathOperator*{\argmax}{arg\,max}

\title{Устаревание драфта EAGLE после дообучения target}
\subtitle{Карта, предсказание по весам, механизм и дешёвая починка без генерации}
\date{15 октября 2026}
\author{}

\begin{document}

\maketitle

\begin{frame}{Проблема}
\small
\begin{itemize}
\item EAGLE-3 ускоряет генерацию в 2--4 раза, но драфт обучают под \alert{конкретный} target.
\item Target постоянно дообучают: LoRA под домены, SFT, RL, дистилляция рассуждений.
\item Драфт остаётся старым. Переобучение (EDA и аналоги) --- генерация новым target и часы обучения.
\end{itemize}
\vspace{4pt}
\begin{block}{Вопросы}
\begin{enumerate}
\item Насколько драфт устаревает на реальных дообучениях?
\item Можно ли это предсказать, не генерируя текст?
\item Почему драфт не подстраивается сам, хотя читает признаки target?
\item Можно ли починить драфт \textbf{без генерации и без обучения}?
\end{enumerate}
\end{block}
\end{frame}

\begin{frame}{Постановка}
\small
Драфт читает три слоя target и сжимает их:
\[
u_t=W_{\fc}\big[h^{(2)}_t;\,h^{(16)}_t;\,h^{(29)}_t\big],\qquad W_{\fc}\in\mathbb R^{4096\times12288}
\]
Дерево 60 токенов глубины 7, $\tau$ --- принятых токенов за проход target.

\vspace{6pt}
Target дообучили: $\theta_0\to\theta_1$, драфт $\phi_0$ прежний. Потеря на области $d$:
\[
L_d(\theta_1)=1-\frac{\tau_d(\phi_0;\theta_1)}{\tau_d(\phi_0;\theta_0)}
\]
\small Протокол: код авторов EAGLE, официальный драфт Llama-3.1-8B-Instruct, 96 запросов
(общие, математика, код, медицина, инструменты, русский), 10 публичных дообучений, Colab A100.
\end{frame}

\begin{frame}{1. Дообучение делает драфт устаревшим}
\begin{columns}[T]
\begin{column}{0.58\textwidth}
\includegraphics[width=\textwidth]{fig/fig_pitch_1_ladder.png}
\end{column}
\begin{column}{0.40\textwidth}
\small
\begin{itemize}
\item LoRA: \textbf{0--10\%}
\item полное по домену: \textbf{10--26\%}, во всех областях
\item рассуждения: \alert{\textbf{40--60\%}}, ускорение $2{,}34\times\to1{,}5\times$
\item свой драфт R1-Distill возвращает всё $\Rightarrow$ дело в несоответствии драфта
\item русский: 56\% токенов вне словаря драфта, $\tau=1{,}67$
\end{itemize}
\end{column}
\end{columns}
\end{frame}

\begin{frame}{2. Устаревание видно без генерации}
\begin{columns}[T]
\begin{column}{0.62\textwidth}
\includegraphics[width=\textwidth]{fig/fig_pitch_2_predictor.png}
\end{column}
\begin{column}{0.36\textwidth}
\small
\[
\rho_W=\operatorname{mean}_{o,\,down}\frac{\|W_1-W_0\|_F}{\|W_0\|_F}
\]
\begin{itemize}
\item только веса: \textbf{r = 0,98}, ρ = 0,87
\item сдвиг состояний: r = 0,93
\item лучше, чем изменение поведения target (0,86)
\item \textit{Оговорка:} держится на тяжёлых изменениях; умеренные не ранжирует
\end{itemize}
\end{column}
\end{columns}
\end{frame}

\begin{frame}{3. Механизм: перенос против искажения}
\small
Делим токены: target \textbf{не изменил} ответ ($S$) или \textbf{изменил} ($C$). Совпадение драфта на первом шаге:
\[
T=a_C(f^1)-a_C(f^0)\ \ (\text{перенос}),\qquad D=a_S(f^0)-a_S(f^1)\ \ (\text{искажение})
\]
\begin{columns}[T]
\begin{column}{0.56\textwidth}
\includegraphics[width=0.92\textwidth]{fig/fig_pitch_3_mechanism.png}
\end{column}
\begin{column}{0.42\textwidth}
\footnotesize
\begin{itemize}
\item перенос есть везде: \pos{+6…+23 п.п.}
\item искажение только у тяжёлых: \negv{−3…−19 п.п.}
\item «драфт сам подстраивается» = $T-D$, у тяжёлых $\approx0$
\item искажение --- \textbf{направление}, не масштаб; 75--80\% через слой 29; бьёт по служебным токенам
\end{itemize}
\end{column}
\end{columns}
\end{frame}

\begin{frame}{4. Где живёт искажение}
\begin{itemize}
\item Поэлементное, полноранговое и низкоранговое восстановление «новые $\to$ старые признаки»: \textbf{30 вариантов}.
\item Лучшее возвращает \textbf{47--57\%} разрыва на первом шаге драфта --- и это 1--4 направления, а не 64.
\item Перенос сохраняется: у Nemotron на $C$ $0{,}281\to0{,}320$ (со старыми признаками $0{,}204$).
\item Вторая половина искажения специфична для токенов и линейно не возвращается.
\end{itemize}
\vspace{6pt}
\begin{alertblock}{Вывод}
Глобальная часть искажения почти одномерна $\Rightarrow$ её можно убрать поправкой ранга 1--2 на слой.
\end{alertblock}
\end{frame}

\begin{frame}{5. Метод: поправка ранга 2, встроенная в $W_{\fc}$}
\small
Для слоя $j$: $P_j=\mathrm{orth}[\mu_j;v_j]$ --- средний сдвиг и главное направление сдвига. Ridge:
\[
\min_{R_j}\big\|(X^0-X^1)P_j^\top-[X^1,\mathbf 1]R_j\big\|_F^2+\lambda\|R_j\|_F^2,\qquad
x'=x(I+A_j)+c_j,\ \ \operatorname{rank}A_j=2
\]
Встраиваем в драфт:
\[
W_{\fc}'=W_{\fc}\,\mathrm{blockdiag}\big(I+A_j^\top\big)_{j=2,16,29},\qquad b_{\fc}=W_{\fc}\,[c_2;c_{16};c_{29}]
\]
\begin{itemize}
\item данные: прямые проходы $\theta_0$ и $\theta_1$ на 24 текстах исходной модели
\item \textbf{новый target ничего не генерирует, драфт не обучается}, ~1 минута на A100
\item вариант зафиксирован \textbf{до} проверки τ
\end{itemize}
\end{frame}

\begin{frame}{6. Результат: τ дерева на ответах самого изменённого target}
\begin{columns}[T]
\begin{column}{0.52\textwidth}
\includegraphics[width=\textwidth]{fig/fig_pitch_4_fix.png}
\end{column}
\begin{column}{0.46\textwidth}
\footnotesize
\begin{tabular}{@{}lcc@{}}
\toprule
 & $\Delta\tau$ & 95\% интервал \\
\midrule
Nemotron-Nano & \pos{\textbf{+8,6\%}} & [+7,4; +10,0] \\
R1-Distill & \pos{\textbf{+9,6\%}} & [+8,4; +10,8] \\
ToolACE (контроль) & +0,2\% & [−0,5; +1,0] \\
abliterated & \negv{−9,7\%} & [−11,0; −8,3] \\
\bottomrule
\end{tabular}

\vspace{6pt}
\begin{itemize}
\item лучше 90--94\% запросов у рассуждающих моделей
\item ускорение ≈ $1{,}51\to1{,}65\times$ и $1{,}74\to1{,}86\times$
\item сравнение на общем начале ответов
\end{itemize}
\end{column}
\end{columns}
\end{frame}

\begin{frame}{Почему не везде: когда включать поправку}
\begin{center}
\begin{tabular}{@{}lccc@{}}
\toprule
 & поворот выхода fc, $\cos(u^0,u^1)$ & масштаб, $\|u^1\|/\|u^0\|$ & эффект \\
\midrule
Nemotron-Nano & 0,75 & 1,49 & \pos{+8,6\%} \\
R1-Distill & 0,82 & 1,13 & \pos{+9,6\%} \\
abliterated & 0,94 & 0,79 & \negv{−9,7\%} \\
ToolACE & 0,99 & 1,05 & +0,2\% \\
\bottomrule
\end{tabular}
\end{center}
\begin{itemize}
\item у рассуждающих моделей искажение --- \textbf{поворот}: на него поправка и рассчитана
\item у abliterated --- \textbf{сжатие}; $u$ идёт в остаточный поток драфта без нормировки
\item гипотеза правила: включать при $\cos(u^0,u^1)<0{,}9$ (4 точки, задним числом)
\end{itemize}
\end{frame}

\begin{frame}{Новизна}
\small
\begin{tabular}{@{}p{0.47\textwidth}p{0.47\textwidth}@{}}
\toprule
Уже известно & Наше \\
\midrule
EDA: адаптация драфта к дообученному target через генерацию и обучение &
починка входа драфта \textbf{из прямых проходов}, без генерации и обучения \\
AngelSpec, Draft-OPD: драфт режима без рассуждений хуже на рассуждениях &
то же для \textbf{изменения весов}: −40…−60\%; систематическая карта на 10 моделях \\
SPEED-Bench, AdaSpec, SpecVocab: урезанный словарь вредит языкам &
\textbf{предсказание устаревания по весам} (r = 0,98) \\
«EAGLE подстраивается сам, т.к. читает признаки» &
разложение \textbf{перенос против искажения}, искажение почти одномерно \\
\bottomrule
\end{tabular}
\end{frame}

\begin{frame}{Выводы}
\begin{enumerate}
\item Дообучение, особенно под рассуждения, \alert{съедает 40--60\%} выгоды EAGLE.
\item Устаревание \textbf{предсказуемо по весам} --- без текстов и прогонов.
\item Новые признаки передают драфту новое поведение, но у тяжёлых изменений \textbf{искажают} его работу; искажение
почти одномерно.
\item Поправка ранга 2 за минуту, без генерации и обучения: \pos{\textbf{+9--10\% τ}} на моделях с рассуждениями.
\item Это «первая помощь»: свой драфт даёт +100\%; поправка не универсальна (abliterated).
\end{enumerate}
\end{frame}

\begin{frame}{Ограничения и следующие шаги}
\begin{columns}[T]
\begin{column}{0.48\textwidth}
\textbf{Ограничения}
\small
\begin{itemize}
\item одна база (Llama-3.1-8B), один драфт
\item 16 запросов на область
\item метод проверен на 4 изменениях; правило --- задним числом
\item ускорение по времени оценено через τ
\end{itemize}
\end{column}
\begin{column}{0.48\textwidth}
\textbf{Дальше}
\small
\begin{itemize}
\item заранее объявленная проверка на 2--3 новых рассуждающих моделях + замер скорости (~1,5 ч A100)
\item вторая семья: Qwen3-8B
\item поправка, учитывающая масштаб
\item направления поправки из весов --- шаг к версии «только по весам»
\end{itemize}
\end{column}
\end{columns}
\end{frame}

\end{document}
